\subsection{The lift of families of elements of an associated graded module}

Let A be a filtered commutative ring, \((\mathbf{A}_{n})_{n\in\mathbf{Z}}\) its filtration and C a subring of \(A_{0}\) such that \(C\cap A_{1}=\{0\}\) . The restriction to C of the canonical mapping \(\mathbf{A}_{0}\to\mathbf{A}_{0}/\mathbf{A}_{1}=\mathrm{gr}_{0}(\mathbf{A})\) is then injective, which allows us to identify C with a submodule of \(\mathrm{gr}_{0}(\mathbf{A})\) ; this is what we shall usually do in similar cases. If \(A_{1}\neq A_{0}\) and K is any subfield of \(A_{0}\) , then \(K\cap A_{1}=\{0\}\) since \(K\cap A_{1}\) is an ideal of K not containing 1; then K may be identified with a subfield of \(\mathrm{gr}_{0}(\mathbf{A})\) .

PROPOSITION 10. Let A be a filtered commutative ring and \((\mathbf{A}_n)\) its filtration; suppose that there exists a subring \(\mathbf{C}\) of \(\mathbf{A}_0\) such that \(\mathbf{C} \cap \mathbf{A}_1 = \{0\}\) and \(\mathbf{C}\) is identified with a subring of \(\operatorname{gr}_0(\mathbf{A})\) . Let \((x_i)_{1 \leqslant i \leqslant q}\) be a finite family of elements of \(\mathbf{A}\) ; suppose that \(x_i \in \mathbf{A}_{n_i}\) for \(1 \leqslant i \leqslant q\) and let \(\xi_i\) be the class of \(x_i\) in \(\operatorname{gr}_{n_i}(\mathbf{A})\) for \(1 \leqslant i \leqslant q\) .

\begin{enumerate}
\def\labelenumi{(\roman{enumi})}
\item
  If the family \((\xi_{i})\) of elements of \(\operatorname{gr}(\mathbf{A})\) is algebraically free over \(\mathbf{C}\) , the family \((x_{i})\) is algebraically free over \(\mathbf{C}\) .
\item
  If the filtration on \(\mathbf{A}\) is exhaustive and discrete and \((\xi_{i})\) is a system of generators of the C-algebra \(\operatorname{gr}(\mathbf{A})\) , then \((x_{i})\) is a system of generators of the C-algebra \(\mathbf{A}\) .
\end{enumerate}

Let \(\mathbf{A}'\) be the polynomial algebra \(\mathbf{C}[\mathbf{X}_1, \ldots, \mathbf{X}_q]\) over \(\mathbf{C}\) ; let \(\mathbf{A}'\) be given the graduation \((\mathbf{A}_n')\) of type \(\mathbf{Z}\) where \(\mathbf{A}_n'\) is the set of C-linear combinations of the monomials \(\mathbf{X}_1^{s(1)} \ldots \mathbf{X}_q^{s(q)}\) such that \(\sum_{i=1}^{q} n_i s(i) = n\) . Let \(u\) be the homomorphism \(f \mapsto f(x_1, \ldots, x_q)\) from the C-algebra \(\mathbf{A}'\) to the C-algebra \(\mathbf{A}\) ; by definition, \(u(\mathbf{A}_n') \subset \mathbf{A}_n\) for all \(n \in \mathbf{Z}\) and hence \(u\) is compatible with the filtrations ( \(\mathbf{A}'\) being given its filtered ring structure associated with its graded ring structure, cf.~no. 1, Example 1). This being so, the hypothesis of (i) means that

\[
\operatorname{gr} (u): \mathrm{A} ^ {\prime} = \operatorname{gr} (\mathrm{A} ^ {\prime}) \rightarrow \operatorname{gr} (\mathrm{A})
\]

is injective; as the filtration on A' is exhaustive and separated, Corollary 1 to Theorem 1 of no. 8 may be applied and u is injective, which proves the conclusion of (i). Similarly, the hypothesis (ii) on the \((\xi_{i})\) means that gr(u) is surjective; as A is discrete and its filtration is exhaustive, Corollary 2 to Theorem 1 of no. 8 may be applied and u is surjective, which proves the conclusion of (ii).

PROPOSITION 11. Let A be a complete Hausdorff filtered commutative ring, C a subring of \(\mathbf{A}_0\) such that \(\mathbf{C} \cap \mathbf{A}_1 = 0\) and \((x_i)_{1 \leqslant i \leqslant q}\) a finite family of elements of A such that \(x_i \in \mathbf{A}_{n_i}\) where \(n_i > 0\) for \(1 \leqslant i \leqslant q\) ; let \(\xi_i\) be the class of \(x_i\) in \(\operatorname{gr}_{n_i}(\mathbf{A})\) for \(1 \leqslant i \leqslant q\) . (i) There exists a unique C-homomorphism \(v\) from the algebra of formal power series \(\mathbf{A}'' = \mathbf{C}[[\mathbf{X}_1, \ldots, \mathbf{X}_q]]\) to A such that \(v(\mathbf{X}_i) = x_i\) for \(1 \leqslant i \leqslant q\) .

\begin{enumerate}
\def\labelenumi{(\roman{enumi})}
\setcounter{enumi}{1}
\item
  If the family \((\xi_{i})\) is algebraically free over \(\mathbf{C}\) , the homomorphism \(v\) is injective.
\item
  If the filtration on \(\mathbf{A}\) is exhaustive and the family \((\xi_{i})\) is a system of generators of the C-algebra \(\operatorname{gr}(\mathbf{A})\) , the homomorphism \(v\) is surjective.
\end{enumerate}

As \(n_i \geqslant 1\) for all \(i, \sum_{i=1}^{q} n_i s(i) \geqslant \sum_{i=1}^{q} s(i)\) for every monomial \(X_1^{s(1)} \ldots X_q^{s(q)}\) and on the other hand \(\sum_{i=1}^{q} n_i s(i) \leqslant r. \sum_{i=1}^{q} s(i)\) if \(r\) is the greatest of the \(n_i\) . If \(A_n''\) denotes the set of formal power series whose non-zero terms \(a_s X_1^{s(1)} \ldots X_q^{s(q)}\) satisfy \(\sum_{i=1}^{q} n_i s(i) \geqslant n\) , it follows from no. 6, Corollary to Proposition 6 that \(A''\) is Hausdorff and complete with the exhaustive filtration \((A_n'')\) and that

\[
\mathrm{A} ^ {\prime} = \mathrm{C} [ \mathrm{X} _ {1}, \dots , \mathrm{X} _ {q} ]
\]

is dense in A''; moreover the homomorphism u defined in the proof of Proposition 10 is continuous on A' and can be extended uniquely to a continuous homomorphism v: A''→A, since A is Hausdorff and complete (General Topology, Chapter III, § 3, no. 3, Proposition 5), which proves (i); also, gr(A'') = gr(A') and gr(v) = gr(u); (ii) and (iii) then follow respectively from Corollaries 1 and 2 to Theorem 1 of no. 8 in view of the hypotheses on A.

The conclusion of (ii) (resp. (iii)) of the proposition is sometimes expressed by saying that the family \((x_{i})\) is formally free over C (resp. a formal system of generators of A).

PROPOSITION 12. Let A be a filtered ring, E a filtered A-module and \((\mathbf{A}_{n})\) and \((\mathbf{E}_{n})\) the respective filtrations on A and E. Suppose that A is complete and the filtration \((\mathbf{E}_{n})\) is exhaustive and separated. Let \((x_{i})_{i\in\mathbf{I}}\) be a finite family of elements of E and for \(i\in I\) let \(n(i)\) be an integer such that \(x_{i}\in\mathbf{E}_{n(i)}\) ; finally let \(\xi_{i}\) be the class of \(x_{i}\) in \(\mathrm{gr}_{n(i)}(\mathbf{E})\) . Then, if \((\xi_{i})\) is a system of generators of the gr(A)-module gr(E), \((x_{i})\) is a system of generators of the A-module E.

In the A-module \(\mathbf{L} = \mathbf{A}_s^{\mathbf{I}}\) let \(\mathbf{L}_n\) denote the set \((a_i)\) such that \(a_i \in \mathbf{A}_{n - n(i)}\) for all \(i \in \mathbf{I}\) ; if \(p\) and \(q\) are the least and greatest of the \(n(i)\) , then \(\mathbf{A}_{n - q}^{\mathbf{I}} \supset \mathbf{L}_n \supset \mathbf{A}_{n - p}^{\mathbf{I}}\) and the topology defined on \(\mathbf{L}\) by the definition \((\mathbf{L}_n)\) is the same as the product topology; hence \(\mathbf{L}\) is a complete filtered A-module. As \(\mathbf{L}\) is free, there exists an

A-linear mapping \(u: \mathbf{L} \to \mathbf{E}\) such that \(u((a_i)) = \sum_{i \in \mathbf{I}} a_i x_i\) and it is obviously compatible with the filtrations; we must prove that \(u\) is surjective and for this it is sufficient, by virtue of Corollary 2 to Theorem 1, no. 8, to show that

\[
\operatorname{gr} (u): \operatorname{gr} (\mathbf {L}) \rightarrow \operatorname{gr} (\mathbf {E})
\]

is surjective or also that, for all \(x \in \mathbf{E}_n\) , there exist a family \((a_i)\) such that \(a_i \in \mathbf{A}_{n-n(i)}\) for all \(i \in I\) and \(x \equiv \sum_{i \in I} a_i x_i \pmod{\mathbf{E}_{n+1}}\) . Let \(\xi\) be the class of \(x\) in \(\mathrm{gr}_n(\mathbf{E})\) ; since the \(\xi_i\) generate the \(\mathrm{gr}(\mathbf{A})\) -module \(\mathrm{gr}(\mathbf{E})\) , there exist \(\alpha_i \in \mathrm{gr}(\mathbf{A})\) such that \(\xi = \sum_{i \in I} \alpha_i \xi_i\) and we may assume that \(\alpha_i \in \mathrm{gr}_{n-n(i)}(\mathbf{A})\) by replacing if need be \(\alpha_i\) by its homogeneous component of degree \(n - n(i)\) . Then \(\alpha_i\) is the image of an element \(a_i \in \mathbf{A}_{n-n(i)}\) and the family \((a_i)\) has the required property.

COROLLARY 1. Let A be a complete filtered ring and E a filtered A-module whose filtration is exhaustive and separated. If gr(E) is a finitely generated (resp. Noetherian) gr(A)-module, then E is a finitely generated (resp. Noetherian) A-module.

If \(\mathrm{gr}(\mathrm{E})\) is finitely generated, there is a finite system of homogeneous generators and Proposition 12 shows that E is finitely generated. Suppose now that \(\mathrm{gr}(\mathrm{E})\) is Noetherian and let F be a submodule of E; the filtration induced on F by that on E is exhaustive and separated and \(\mathrm{gr}(\mathrm{F})\) is identified with a sub-gr(A)-module of gr(E) (no. 4, Proposition 2) and hence is finitely generated by hypothesis; we conclude that F is a finitely generated A-module and hence E is Noetherian.

COROLLARY 2. Let A be a complete Hausdorff filtered ring whose filtration is exhaustive. If gr(A) is a left Noetherian ring, so is A.

It is sufficient to apply Corollary 1 with \(\mathbf{E} = \mathbf{A}_s\) .

COROLLARY 3. Let A be a complete filtered ring, \((\mathbf{A}_{n})\) its filtration, E a Hausdorff filtered A-module, \((\mathbf{E}_{n})\) its filtration and F a finitely generated submodule of E; suppose that \(A_{0} = A\) and \(E_{0} = E\) .

\begin{enumerate}
\def\labelenumi{(\roman{enumi})}
\item
  If, for all \(k \geqslant 0\) , \(\mathbf{E}_k = \mathbf{E}_{k+1} + \mathbf{A}_k\mathbf{F}\) , then \(\mathbf{F} = \mathbf{E}\) .
\item
  If it is further supposed that the filtration on \(\mathbf{E}\) is derived from that on \(\mathbf{A}\) (no. 1, Example 2), the relation \(\mathbf{E} = \mathbf{E}_1 + \mathbf{F}\) implies \(\mathbf{F} = \mathbf{E}\) .
\end{enumerate}

Let \(\xi_{i}(1\leqslant i\leqslant n)\) be the classes mod. \(\mathbf{E}_1\) of a finite system of generators of \(\mathbf{F}\) . It follows from the given hypothesis that for all \(k\geqslant 0\) every element of \(\mathrm{gr}_k(\mathrm{E})\)

can be expressed in the form \(\sum_{i=1}^{\infty} \alpha_i \xi_i\) where \(\alpha_i \in \mathrm{gr}(A)\) ; the \(\xi_i\) therefore generate the \(\mathrm{gr}(A)\) -module \(\mathrm{gr}(E)\) , which proves (i) by virtue of Proposition 12. If the filtration on \(E\) is derived from that on \(A\) , the relation \(E = E_1 + F\) implies

\[
\begin{array}{r} \mathbf {E} _ {k} = \mathbf {A} _ {k} \mathbf {E} = \mathbf {A} _ {k} \mathbf {E} _ {1} + \mathbf {A} _ {k} \mathbf {F} = \mathbf {A} _ {k} \mathbf {A} _ {1} \mathbf {E} + \mathbf {A} _ {k} \mathbf {F} \subset \mathbf {A} _ {k + 1} \mathbf {E} + \mathbf {A} _ {k} \mathbf {F} \\ = \mathbf {E} _ {k + 1} + \mathbf {A} _ {k} \mathbf {F} \subset \mathbf {E} _ {k}, \end{array}
\]

whence (ii).

PROPOSITION 13. Let A be a ring, m a two-sided ideal of A contained in the Jacobson radical of A and E an A-module. Let A and E be given the m-adic filtrations (no. 1, Example 3). Suppose that one of the following conditions holds:

\begin{enumerate}
\def\labelenumi{(\alph{enumi})}
\item
  E is a finitely generated A-module and A is Hausdorff;
\item
  m is nilpotent.
\end{enumerate}

For E to be a free A-module, it is necessary and sufficient that E/mE be a free (A/m)-module and that E satisfy property (GR) (no. 8).

If E is a free A-module and \((e_{\lambda})\) a basis of E, \(m^{k}E\) is the direct sum of the submodules \(m^{k}e_{\lambda}\) of E for all \(k \geqslant 0\) (Algebra, Chapter II, §3, no. 7, Remark); then \(m^{k}E/m^{k+1}E\) is identified with the direct sum of the \(m^{k}e_{\lambda}/m^{k+1}e_{\lambda}\) (Algebra, Chapter II, §1, no. 6, Proposition 7). We deduce first (for k = 0) that the classes \(1 \otimes e_{\lambda}\) of the \(e_{\lambda}\) in E/mE = (A/m) \(\otimes_{A} E\) form a basis of the (A/m)-module E/mE, since the canonical mapping

\[
\left(\mathfrak {m} ^ {k} / \mathfrak {m} ^ {k + 1}\right) \otimes_ {\mathbf {A}} (\mathrm{E} / \mathfrak {m} \mathrm{E}) \rightarrow \mathfrak {m} ^ {k} \mathrm{E} / \mathfrak {m} ^ {k + 1} \mathrm{E}
\]

is bijective for all \(k \geqslant 0\) ; hence \(\gamma_{\mathbf{E}}\) is bijective. Note that this part of the proof uses neither condition (a) nor condition (b).

Suppose conversely that the conditions of the statement hold and let \((x_{\iota})_{\iota \in I}\) be a family of elements of E whose classes mod. mE form a basis of the (A/m)-module E/mE; let L be the free A-module \(\mathrm{A}_s^{(I)}\) , \((f_{\iota})_{\iota \in I}\) its canonical basis and \(u: \mathbf{L} \to \mathbf{E}\) the A-linear mapping such that \(u(f_{\iota}) = x_{\iota}\) for all \(\iota \in I\) . The hypotheses already imply that \(u\) is surjective (Chapter II, §3, no. 2, Corollary 1 to Proposition 4) and it remains to prove that \(u\) is injective. Now, each of the hypotheses (a) and (b) implies that A is Hausdorff and hence so is L with the m-adic filtration, since \(\mathfrak{m}^k\mathbf{L} = (\mathfrak{m}^k)^{(I)}\) (Algebra, Chapter II, §3, no. 7, Remark) and \(\operatorname{gr}(\mathbf{L})\) is identified with \(\operatorname{gr}(\mathbf{A}) \otimes_{\mathbf{A} / \mathfrak{m}} (\mathbf{L} / \mathfrak{m}\mathbf{L})\) from the first part of the proof; the homomorphism \(u\) is compatible with the filtrations and it is possible to write \(\operatorname{gr}(u) = \gamma_{\mathbb{E}} \circ v\) where \(v\) is the bijection of \(\operatorname{gr}(\mathbf{L})\) onto

\[
\operatorname{gr} (\mathbf {A}) \otimes_ {\mathbf {A} / \mathfrak {m}} (\mathbf {E} / \mathfrak {m} \mathbf {E})
\]

mapping the class of \(f_{t}\) mod. mM onto \(1 \otimes \bar{x}_{t}\) , where \(\bar{x}_{t}\) is the class of \(x_{t}\) mod. mE. The hypothesis then implies that \(\operatorname{gr}(u)\) is injective and the conclusion follows with the aid of Corollary 1 to Theorem 1, no. 8.

